\chapter{Power Grid Reconstruction Methodologies}

\section{Xiong et al. -- OSM/PyPSA}

\begin{description}

    \item[\textbf{Reference:}]
    Xiong et al., ``Modelling the high-voltage grid using open data for
    Europe and beyond''.

    \item[\textbf{Objective:}]
    Develop an open and reproducible methodology for constructing a
    topologically connected and simulation-ready representation of the
    European high-voltage transmission grid using primarily OpenStreetMap
    data. The methodology is designed to overcome the lack of openly
    available detailed transmission-system models and is integrated into
    the PyPSA-Eur workflow.

    \item[\textbf{Input data:}] \hfill
    \begin{itemize}
        \item OpenStreetMap AC transmission lines and cables.
        \item OpenStreetMap substations.
        \item OpenStreetMap HVDC connections.
        \item Geographical coordinates and geometries.
        \item Voltage levels.
        \item Number of circuits and cables when available.
        \item Line frequency information.
        \item Standard electrical line types from power-system modelling
        libraries.
    \end{itemize}

    \item[\textbf{Methodology:}] \hfill
    \begin{itemize}
        \item Retrieval of power-system infrastructure from OpenStreetMap
        using the Overpass API.
        \item Extraction of AC lines, cables, substations and HVDC
        infrastructure.
        \item Cleaning and normalisation of incomplete or inconsistent OSM
        attributes.
        \item Selection of high-voltage substations and transmission
        infrastructure.
        \item Identification of line endpoints and conversion of geographical
        infrastructure into network buses and branches.
        \item Detection of transmission lines passing through substations and
        splitting of lines where necessary.
        \item Spatial clustering of nearby buses belonging to the same
        substation.
        \item Connection of different voltage levels within substations using
        transformers.
        \item Reconstruction of a topologically connected transmission
        network.
        \item Separate processing and integration of HVDC connections.
        \item Assignment of standard electrical line types according to
        voltage and available OSM characteristics.
        \item Estimation of missing electrical parameters and transmission
        capacities from standard line types, line length and circuit
        information.
        \item Conversion of the resulting network into a representation
        directly usable by PyPSA and PyPSA-Eur.
    \end{itemize}

    \item[\textbf{Inferred / generated information:}] \hfill
    \begin{itemize}
        \item Network buses.
        \item Connected AC transmission lines and cables.
        \item HVDC links.
        \item Transformers between voltage levels.
        \item Connected high-voltage network topology.
        \item Estimated resistance and reactance.
        \item Estimated transmission capacities.
        \item Electrical connections not explicitly represented in the
        original geographical data.
        \item Simulation-ready PyPSA network representation.
    \end{itemize}

    \item[\textbf{Output:}]
    An open European high-voltage transmission network containing buses,
    AC lines and cables, HVDC links and transformers. The processed network
    can be exported as tabular datasets such as CSV files and used directly
    within the PyPSA/PyPSA-Eur modelling ecosystem.

    \item[\textbf{Validation:}]
    The reconstructed network is compared with ENTSO-E transmission-network
    statistics and other available official data. The study reports strong
    agreement in total transmission-route and circuit lengths across European
    countries. Electrical line parameters are additionally compared with an
    official German transmission-system model, and the reconstructed network
    is tested in a full PyPSA-Eur optimisation covering 8760 hours. The study
    also investigates discrepancies in Spain and compares the reconstructed
    topology with official Spanish geographical information.

    \item[\textbf{Main assumptions:}] \hfill
    \begin{itemize}
        \item Missing circuit information may be completed using predefined
        assumptions.
        \item Lines are assigned standard electrical types according to
        voltage and available characteristics.
        \item Electrical parameters such as resistance and reactance are
        derived from standard line types when official parameters are not
        available.
        \item Transmission capacities are estimated from the assigned
        technical characteristics rather than obtained directly from the TSO.
        \item Nearby buses can be clustered when they are considered to
        represent the same physical substation.
        \item Transformers can be synthetically introduced when buses at
        different voltage levels belong to the same substation.
    \end{itemize}

    \item[\textbf{Main limitations:}] \hfill
    \begin{itemize}
        \item The reconstructed grid is not an exact representation of the
        official transmission network operated by each TSO.
        \item Results depend on the completeness and quality of OpenStreetMap
        data.
        \item Some electrical parameters are inferred from standard equipment
        types rather than obtained from official network data.
        \item Transmission capacities are estimates and should not
        automatically be interpreted as official operational limits.
        \item Some transformers are synthetic because sufficiently complete
        transformer information is not available from OSM.
        \item The methodology mainly focuses on high-voltage transmission
        infrastructure of 200~kV and above.
    \end{itemize}

    \item[\textbf{Application in EAGLE:}]
    This methodology is particularly relevant for constructing the base
    electrical topology of the EAGLE Spanish transmission-network model.
    Its workflow can be used to transform raw OpenStreetMap geographical
    information into connected buses, transmission lines, transformers and
    HVDC links. The proposed use of standard line types provides a practical
    method for estimating missing resistance, reactance and transmission
    capacities when official values are unavailable. In addition, the
    resulting network is directly compatible with PyPSA/PyPSA-Eur, making
    the methodology especially suitable if PyPSA is selected as one of the
    EAGLE simulation and optimisation frameworks.

\end{description}


\section{Rivera et al. -- OSM/Transnet}

\begin{description}

    \item[\textbf{Reference:}]
    Rivera et al., ``Automatic Generation of Real Power Transmission Grid
    Models From Crowdsourced Data'', IEEE Transactions on Smart Grid, 2019.

    \item[\textbf{Objective:}]
    Automatically reconstruct transmission grid models from incomplete and
    unstructured crowdsourced OpenStreetMap (OSM) data. The methodology aims
    to infer electrical connections between power lines, substations and
    generators and transform the resulting topology into a simulation-ready
    electrical network model.

    \item[\textbf{Input data:}] \hfill
    \begin{itemize}
        \item OpenStreetMap power lines.
        \item Substations.
        \item Power plants and generators.
        \item Geographical coordinates and geometries.
        \item Voltage levels and other available OSM attributes.
        \item Population data for load estimation.
        \item Average electricity consumption data for load estimation.
    \end{itemize}

    \item[\textbf{Methodology:}] \hfill
    \begin{itemize}
        \item Extraction of power-related objects from OpenStreetMap.
        \item Identification of relevant transmission voltage levels.
        \item Filtering of transmission lines and possible network endpoints
        for each voltage level.
        \item Spatial identification of intersections between power lines,
        substations, power plants and generators.
        \item Recursive traversal of connected lines to identify the endpoints
        of electrical circuits.
        \item Automatic inference of power circuit relations that are not
        explicitly defined in OpenStreetMap.
        \item Treatment of T-junctions by dividing them into individual
        electrical relations.
        \item Construction of a connected transmission network topology.
        \item Estimation of missing electrical parameters.
        \item Estimation of nodal loads using geographical population
        information and electricity consumption.
        \item Conversion of the reconstructed network into Common Information
        Model (CIM).
        \item Conversion of CIM models into MATLAB/Simulink models for
        power-flow simulation.
    \end{itemize}

    \item[\textbf{Inferred / generated information:}] \hfill
    \begin{itemize}
        \item Connected transmission network topology.
        \item Electrical circuit relations between network elements.
        \item Connections between lines, substations and generators.
        \item Estimated line resistance and reactance.
        \item Estimated nodal loads.
        \item CIM representation of the reconstructed network.
        \item Simulation-ready MATLAB/Simulink network model.
    \end{itemize}

    \item[\textbf{Output:}]
    Reconstructed transmission network models that can be stored in JSON
    and CSV formats and subsequently converted to CIM and MATLAB/Simulink
    representations suitable for power-flow simulations.

    \item[\textbf{Validation:}]
    The methodology was evaluated for 14 countries by comparing the
    automatically inferred grid with manually defined OSM circuit relations
    and official transmission-network statistics. The topology-inference
    algorithm achieved high recall when compared with manually defined OSM
    relations. For Spain, the reconstructed 220--400~kV network represented
    approximately 63\% of the officially reported transmission-line length
    using the OSM data considered in the study.

    \item[\textbf{Main assumptions:}] \hfill
    \begin{itemize}
        \item Only transmission voltage levels with sufficient representation
        in OSM are considered.
        \item Missing conductor characteristics are inferred from technical
        assumptions associated with voltage levels.
        \item Line resistance and reactance are approximated from line length
        and assumed conductor characteristics.
        \item Electricity demand is geographically distributed using
        population information and average electricity consumption.
        \item The network is assumed to be a balanced three-phase system for
        the demonstrated power-flow simulations.
    \end{itemize}

    \item[\textbf{Main limitations:}] \hfill
    \begin{itemize}
        \item The quality of the reconstructed network depends strongly on
        the completeness and accuracy of OpenStreetMap.
        \item The Spanish evaluation is based on old OSM data and should not
        be interpreted as a current representation of the Spanish grid.
        \item The reported 63\% value for Spain is primarily a comparison of
        transmission-line length and does not represent 63\% accuracy of all
        electrical parameters or network topology.
        \item Resistance and reactance are rough estimates rather than
        official transmission-system parameters.
        \item Nodal loads are estimated from population data rather than
        obtained from actual electrical measurements.
        \item Intentional network loops can be ignored by the proposed
        inference algorithm to reduce computational complexity.
    \end{itemize}

    \item[\textbf{Application in EAGLE:}]
    The main value for EAGLE is methodological rather than as a direct source
    of Spanish transmission-grid data. The proposed spatial connectivity
    inference can be adapted to reconstruct electrical relations when
    OpenStreetMap data do not explicitly define them. Its approaches for
    estimating missing network parameters and nodal loads also provide
    possible strategies for completing incomplete EAGLE datasets. The CIM
    conversion demonstrates an additional method for transforming the
    reconstructed topology into a standard electrical network representation.

\end{description}
